% from ega0-19-fr.tex
\oldpage[0\textsubscript{IV-1}]{88}

  appliquant (i) à un cocycle définissant cette extension (18.4.3), on en
  conclut qu'il y a un idéal ouvert
  $\mathfrak R''\subset\mathfrak R'$ de $B$ tel que l'image réciproque
  $E''$ de $E$ par $B/\mathfrak R''\to B/\mathfrak R$ soit
  $A$-triviale, C.Q.F.D.
\end{enumerate}
\end{proof}

\begin{corollary}[19.4.6]
\label{0.19.4.6-fr}
Soient $A$ un anneau topologique, $B$ une $A$-algèbre topologique qui soit un
$A$-module formellement projectif. Soit $A'$ une $A$-algèbre \emph{munie de
la topologie déduite de celle de $A$}. Supposons en outre que $A'$ soit un
$A$-module fidèlement plat, et que l'une des conditions suivantes soit
remplie~:
\begin{enumerate}
  \item[$1^\circ$] Il existe un système fondamental
  $(\mathfrak J_\lambda)$ d'idéaux ouverts de $A$ et un système fondamental
  $(\mathfrak M_\lambda)$ d'idéaux ouverts de $B$, ayant même ensemble
  d'indices et tels que, pour tout $\lambda$, on ait
  $\mathfrak J_\lambda B\subset\mathfrak M_\lambda$ et que
  $B/\mathfrak M_\lambda$ soit un
  $(A/\mathfrak J_\lambda)$-module projectif de type fini.
  \item[$2^\circ$] $A'$ est un $A$-module projectif de type fini.
\end{enumerate}
Alors, pour que $B'=B\otimes_A A'$ (muni de la topologie produit tensoriel)
soit une $A'$-algèbre formellement lisse, il faut et il suffit que $B$ soit une
$A$-algèbre formellement lisse.
\end{corollary}

\begin{proof}
La suffisance de la condition est contenue dans (19.3.5), (iii), sans aucune
hypothèse supplémentaire sur $B$ ni $A'$. Pour démontrer la réciproque, nous
allons appliquer le critère (19.4.5)~; sous l'hypothèse $2^\circ$, nous
désignerons encore par $(\mathfrak M_\lambda)$ un système fondamental
d'idéaux ouverts de $B$, et, pour tout $\lambda$, par
$\mathfrak J_\lambda$ un idéal ouvert de $A$ tel que
$\mathfrak J_\lambda B\subset\mathfrak M_\lambda$~; dans les deux
cas, nous poserons
$A_\lambda=A/\mathfrak J_\lambda$,
$B_\lambda=B/\mathfrak M_\lambda$,
$A'_\lambda=A_\lambda\otimes_A A'$,
$B'_\lambda=B_\lambda\otimes_A A'$. Soit $f$ un $2$-cocycle de
Hochschild $A_\lambda$-bilinéaire symétrique de
$B_\lambda\times B_\lambda$ dans un $B_\lambda$-module $L$~; par
extension des scalaires, on en déduit un $2$-cocycle de Hochschild
$f'=f\otimes1$, $A'_\lambda$-bilinéaire symétrique, de
$B'_\lambda\times B'_\lambda$ dans $L'=L\otimes_A A'$. Comme par
hypothèse $B'$ est une $A'$-algèbre formellement lisse, il existe, par (19.4.5),
(i), un indice $\mu$ tel que
$\mathfrak M_\mu\subset\mathfrak M_\lambda$, et tel que, si
$\varphi':B'_\mu\to B'_\lambda$ est l'application canonique,
$f'\circ(\varphi'\times\varphi')$ soit un $2$-cobord de Hochschild de
$B'_\mu\times B'_\mu$ dans $L'$~; autrement dit, son image $c'$ dans le
groupe de Hochschild
$\mathrm H^2_{A'_\mu}(B'_\mu,L')^s$ est nulle~; il est clair que si
$\varphi:B_\mu\to B_\lambda$ est l'homomorphisme canonique, et $c$ la
classe de $f\circ(\varphi\times\varphi)$ dans le groupe de Hochschild
$\mathrm H^2_{A_\mu}(B_\mu,L)^s$, $c'$ est l'image canonique de $c$.
Or, si $P_\bullet$ est le complexe relatif aux anneaux $A_\mu$ et
$B_\mu$ défini dans (18.4.5), servant au calcul de
$\mathrm H^2_{A_\mu}(B_\mu,L)^s$, le complexe analogue relatif aux
anneaux $A'_\mu$ et $B'_\mu$ est évidemment
$P_\bullet\otimes_A A'$~; dans l'hypothèse $1^\circ$, la construction de
$P_\bullet$ montre que c'est un $A_\mu$-module projectif de type fini.
On conclut donc de Bourbaki, \emph{Alg.}, chap.~II, 3\textsuperscript{e} éd., §~5,
n\textsuperscript{o}~3, prop.~7 que, sous les deux hypothèses, on a
\[
  \operatorname{Hom}_{A'_\mu}(P_\bullet\otimes_A A',L\otimes_A A')
  =\operatorname{Hom}_{A_\mu}(P_\bullet,L)\otimes_A A'
\]
à un isomorphisme canonique près~; comme $A'$ est un $A$-module plat, on a par
suite (18.4.5)
\[
  \mathrm H^2_{A'_\mu}(B'_\mu,L')^s
  =\mathrm H^2_{A_\mu}(B_\mu,L)^s\otimes_A A'
\]
et l'on peut donc écrire $c'=c\otimes1$~; mais comme $A'$ est un $A$-module
fidèlement plat, l'hypothèse $c'=0$ entraîne $c=0$, ce qui achève la
démonstration.
\end{proof}

\begin{proposition}[19.4.7]
\label{0.19.4.7-fr}
Soient $A$ un anneau topologique préadmissible, $\mathfrak J$ un idéal de
définition de $A$ ($\mathbf 0_{\rm I}$, 7.1.2), $B$ une $A$-algèbre
topologique qui soit un $A$-module formellement projectif. Considérons les
anneaux topologiques quotients $A_0=A/\mathfrak J$,
$B_0=B/\mathfrak JB$~; alors, pour que $B$ soit une $A$-algèbre formellement
lisse, il faut et il suffit que $B_0$ soit une $A_0$-algèbre formellement
lisse.
\end{proposition}

